\chapter[Predicting Failures in BSS Maintenance Operations]{Predicting Failures in \ac{bss} Maintenance Operations}
\label{ch:failure_forecasting_bicing}

Predicting component failures and scheduling maintenance effectively has become a central challenge in the operation of \acp{bss}. The scale of fleets and the importance of reliability increase the need for anticipatory models of when components are likely to fail, as unplanned failures directly affect vehicle availability and user experience. Bikes with worn or failed components cannot be rented or are unsafe, thereby constraining supply, undermining user trust, and ultimately increasing reactive repair costs and downtime.

The complexity of this challenge stems from two main factors: the variability of usage and wear across bikes and components, and the censoring inherent in maintenance records (many components remain in service at the end of the observation period). This Chapter establishes the theoretical and methodological foundations for predicting component-level failures in \acp{bss}, and presents the \ac{pm} framework that combines \ac{sa} models with interpretability techniques for brake pads, wheel spokes, and chains in both \acp{m-b} and \acp{e-b}. The material in this Chapter is adapted from Grau-Escolano et al.~\cite{Grau-Escolano2024}, published in \textit{EPJ Data Science}.

\section{Introduction}

\acp{bss} are a crucial part of urban mobility solutions, providing an eco-friendly alternative to private vehicles. Users notably reduce their reliance on other transportation modes~\cite{Ma2020}, leading to increased physical activity and travel time savings~\cite{Qiu2018}. Additionally, urban environments benefit from reduced fuel consumption and improved economic development~\cite{Qiu2018}. Moreover, \acp{bss} have experienced substantial growth since the 2000s, with 2022's global landscape encompassing nearly 2,000 \acp{bss} and an impressive fleet of approximately 9 million bicycles~\cite{MeddinWorldMapReport2022}. The scale of these systems and technological advancements have enabled the generation of large amounts of data, providing researchers and decision-makers with valuable insights to improve existing systems.

\ac{bss} management is often discussed in terms of mobility and rebalancing; however, maintenance is also important for system reliability and user experience. While the importance of maintenance as a critical activity for companies has been underscored in various studies~\cite{Cho1991, Reinertsen1996}, \ac{bss} maintenance practices have predominantly relied on the traditional approach of applying corrective measures (i.e., addressing issues as they arise by replacing worn-out or damaged parts). Therefore, enhancing the \ac{bss} user experience and operational efficiency could be achieved by developing a \ac{pm} strategy capable of predicting bike component failures before they occur, allowing for timely scheduling of maintenance activities. This approach would not only improve the availability and security of the bike fleet but also reduce costs.

To address this challenge, this Chapter leverages both trip and maintenance datasets provided by \textit{Bicing}. Our research revolves around three key questions. First, we explore bike component failure patterns, seeking to understand how the bike model factor contributes to the wear and tear of distinct parts. Second, we aim to assess the potential of forecasting bike component failures by evaluating the accuracy of various survival models. Specifically, our \ac{pm} system focuses on three key components: brake pads, wheel spokes, and chains, chosen for their substantial representation in the data. Third, we aim to identify the critical factors that influence predicting the longevity of the bike components.


\section{Literature review}

Since the 1990s, the scientific literature has emphasized the importance of maintenance as a critical activity for companies to improve reliability and reduce costs~\cite{Cho1991, Reinertsen1996}. Over time, knowledge and techniques have evolved, transitioning from a corrective maintenance, which involved addressing failures after they occurred, to a preventive maintenance~\cite{Ran2019}, which implements scheduled maintenance based on time intervals to prevent breakdowns. While this approach helps to mitigate most failures, it comes with the drawback of high prevention costs. With the increased computational power, the advancement of \ac{ai}, and the rise of the IoT, a new strategy called \ac{pm} has emerged, which can predict failures before they happen, allowing an even lower failure rate and reduced costs. One of its main challenges is fault prognosis, which focuses on forecasting when a failure is likely to occur~\cite{Wang2017_paradigm}. By accurately predicting failures, maintenance activities can be scheduled in advance, minimizing downtime, and optimizing resource allocation.

Reliability theory plays a significant role in \ac{pm} modeling, and it is essentially the same as \ac{sa}~\cite{Yang2022_prognostic}. \ac{sa} was originally developed in the field of biomedical sciences to examine life tables~\cite{Cox1972}, but its concept of events can be applied to various domains, such as machine component failures. A key challenge in \ac{sa} studies is censoring, which refers to missing data when an event is not observed~\cite{Leung1997, Gijbels2010}. Censoring occurs not due to technical failures but rather due to the nature of the studied event. For instance, if a participant in a clinical trial decides to discontinue his participation before the event of interest occurs. It is precisely the presence of censored data that makes impractical the application of predictive algorithms using the usual statistical and \ac{ml} approaches.

According to~\cite{Wang2019_ml_review}, survival methods can be classified into statistical and \ac{ml} methods. Statistical methods focus on characterizing the distribution of event times and the statistical properties of parameter estimation, such as estimating survival curves. These first models can be further divided into: (1) Non-parametric models (Kaplan-Meier, Nelson-Aalen, Life-Table), which make no assumptions about the underlying distribution and estimate the survival curves directly from the data. (2) Semi-parametric models (Cox model, CoxBoost, Time-dependent Cox), which incorporate both non-parametric estimation of the baseline survival function and parametric estimation of the effects of covariates. (3) Parametric models (Penalized regression, Accelerated Failure Time), which assume a specific distribution for the survival time and estimate the parameters of that distribution. \ac{ml} methods combine traditional \ac{sa} techniques with \ac{ml} algorithms, such as survival trees, Bayesian methods, neural networks, or support vector machines. Advanced \ac{ml} techniques, including ensemble learning, active learning, transfer learning, and multitask learning, have also been applied in the field of \ac{sa}.

These methodologies have been applied in a wide range of fields, including healthcare~\cite{Reddy2015}, reliability~\cite{Modarres2016}, crowdfunding~\cite{Li2016}, student retention~\cite{Ameri2016}, customer lifetime~\cite{Furrer2022}, and unemployment duration analysis~\cite{Kiefer1988}. However, to the best of our knowledge, there are no applications of \ac{sa} to \acp{bss}. Apart from \ac{sa}, scientific literature on bike's \ac{pm} is scarce.~\cite{MountainBike2019} proposed using smartphone vibration readings and support vector machine models to predict the health of mountain bike's components (the rotor, the chain, the wheel bearings, the steering head, and the derailleur cog). However, the scope of this study was not a fleet of bikes but a single bicycle. On the contrary,~\cite{Predictivemaintenancebike2018}'s main objective was to study the cyclists' behavioral patterns in the \ac{bss} of Oslo, Norway, and successfully identify the need for bike maintenance. In this case, random forest models were applied to the ride (destination, duration, and date) and the cyclist (gender and year of birth) data. Finally,~\cite{Matkovic2021} focuses on predicting brakes' performance with \ac{knn}, \ac{lstm} and \ac{xgb} classifiers, using as input physical influences and acceleration/deceleration forces coming from hall and inertial \ac{iot} sensors.
\section{Methodology}
\label{sec:failure_forecasting_methodology}

    \subsection{Data}

\textit{Bicing} has made available for this use case two data sets (Appendix~\ref{app_paper3:section_data_exploration}):

\begin{itemize}
    \item The trips dataset, described in detail in Chapter~\ref{ch:bicing_mobility_patterns}, is used in its entirety for this Chapter, covering the period from April 2019 to December 2022. The same trip filtering as in that Chapter is applied, retaining only trips with duration between 2 and 60 minutes.

    \item The maintenance data set, which comprises a total of 310,000 \acp{mo} corresponding to various bicycle repairs types executed between September 1st, 2020, and January 1st, 2023 (i.e., 2 years and 4 months). Each \ac{mo} record provides the following information: \ac{mo} identifier, date, category, subcategory, bike identifier, and bike model (\ac{m-b} or \ac{e-b}).
    There exist a total of 12 categories and 87 subcategories, encompassing a wide array of actions such as cleaning, greasing, adjusting, or changing bike parts. Categories range from fewer than 100 interventions to 120,000, with the majority of them falling into the brake and wheel categories (66\%). Likewise, subcategories also reveal a wide range of counts, with only 10\% of the repair typologies surpassing 10,000 \acp{mo}.
  \end{itemize}

    \subsection[MOs processing]{\acp{mo} processing} 
    \label{section_mos_processing}

First, the target bike parts for this study needed to be selected. To ensure complete objectivity, subjective \ac{mo} types such as cleaning or greasing were excluded, and the focus was placed on the replacement of bike parts. Also, failures related to wheel tubes were excluded due to their high level of randomness.

\ac{mo} data records specific dates for bike part repairs, however, this data structure is not optimal for conducting \ac{sa}, which typically models the time to an event. To facilitate the use of \ac{sa}, \acp{mo} were transformed into \ac{mo} units. These units represent the time elapsed between two consecutive repairs for the same bike and bike component, and thus, it is the time period in which the bike part under study is operational. Consequently, when working with \ac{mo} units, bike part survival information at the beginning and end of the data set could be lost for each bike. To prevent this information loss, the time from the start of the data set to the first \ac{mo} is considered as one \ac{mo} unit, and the time from the last repair of the bike to the end of the data set is regarded as another \ac{mo} unit. It's worth noting that even though \ac{mo} units may originate from the same bike, each one has been treated as an independent entity.

One key challenge in predicting the occurrence of an event is the presence of censored data, which refers to incomplete information about survival times~\cite{Leung1997, Gijbels2010}. Distinct types of censoring exist, including right-censoring (the event has not occurred by the end of the study), left-censoring (the event occurred before the study started), and interval-censoring (the event occurred within a specific time interval). \textit{Bicing} maintenance data contains left- and right-censored data, since the first and last \ac{mo} units of each bike and repair typology are incomplete. In the first unit, it is not possible to know when this bike part started working, and in the last, when this bike part finally will break. Since \ac{sa} can effectively utilize uncensored and right-censored data to estimate the survival curves and generate predictions, only the left-censored \ac{mo} units were discarded in the training and prediction phases. Specifically, this involves excluding the initial \ac{mo} unit of each bike for the target bike part. Additionally, units without trips and the ones in which \ac{m-b} were upgraded to \ac{e-b} were also discarded.

The data sets used for the survival models follow a format where each row encapsulates all the pertinent information about one subject. Within each entry, there are the details about the subject's survival duration, event occurrence, and aggregated covariate values crucial for the survival function estimation. \ac{mo} units, with their defined start and end dates, facilitate the calculation of covariates that describe the utilization patterns of a bike part and its surrounding environmental factors. In this way, various covariates domains have been incorporated to the model inputs:

\begin{itemize}
\item \textbf{Weather}: weather data was analyzed to incorporate the surrounding environmental conditions of the \ac{mo} unit. This involved considering the daily average temperature (in ºC), total daily precipitation (in mm), average wind direction (in degrees), mean wind speed (in km/h), and average atmospheric pressure (in hPa). After exploring this data and its connection to bike part failures, it was determined that the primary influential factors were the daily average temperature and the mean atmospheric pressure.
\item \textbf{Bike usage}: various metrics were computed, including daily distance traveled (in meters), daily positive and negative inclinations (in meters), and the mean daily speed (in km/h). These calculations were based on the previously collected trip routes and inclinations between each couple of stations. Next, their cumulative versions were examined, and strong correlations were identified. As a result, only two variables were kept: the cumulative daily distance and the mean speed of the \ac{mo} unit.
\item \textbf{Bike model}: due to potential variations in survival curves across the two bicycle models, a binary variable was introduced to distinguish between electric (1) and mechanical (0) bike models.
\item \textbf{Count of repairs for other bike parts during the target \ac{mo} unit}: when analyzing a specific bike part, it becomes pertinent to evaluate the frequency of replacements for another bike part. This approach allows to gain valuable insights; for instance, understanding the number of wheel tubes replaced could aid in predicting potential replacements for components such as tires or wheel rims. Following an examination of the correlation and the variation inflation factor of these repair counts, the following \ac{mo} subcategories were selected: brake tension adjustment, the replacement of the front and rear wheel tubes, and the replacement of the front wheel cover.
\end{itemize}

    \subsection{Models}
To predict the survival time of the bike components, the following statistical and \ac{ml} models have been employed:

\begin{enumerate}
\item \textbf{\acl{cph}} (\acs{cph})\acused{cph} model~\cite{Cox1972} (lifelines' implementation~\cite{Lifelines2021}): is a semi-parametric approach that enables to evaluate how covariates influence the hazard rate of an event as time progresses. This model relies on an underlying assumption known as the proportional hazard assumption, which asserts that the relative risk between two distinct groups remains constant over time. Meeting this assumption simplifies the analytical process and enhances the interpretability of outcomes. Although evaluating the assumption holds theoretical importance for attaining a meaningful interpretation of covariates, adherence to this assumption might not be imperative.~\cite{Stensrud2020} noted that when dealing with a sufficiently large sample size, even minor deviations from the assumption may show up. Furthermore, when the main goal is survival prediction, there is no need to test the proportional hazard assumption, since the main objective is to maximize an score~\cite{lifelines_ph_assumption}. Consequently, for the purpose of this study, the analysis will prioritize predictive accuracy over strict adherence to the assumption.

\item \textbf{\acl{mtlr}} (\acs{mtlr})\acused{mtlr} model~\cite{Yu2021} (pysurvival's implementation~\cite{pysurvival}): stands as an alternative to \ac{cph} when the assumption of proportional hazards does not hold. \ac{mtlr} model relies on a sequence of logistic regression models constructed across distinct time intervals. This allows the estimation of the probability associated with the occurrence of the event of interest within each interval. Consequently, the initial step involves specifying the desired number of intervals, with the present use case opting for the number of days in the maintenance data.

\item \textbf{\acl{csf}} (\acs{csf})\acused{csf} model~\cite{Wright2017} (pysurvival's implementation~\cite{pysurvival}): is a \ac{ml} model that extends Random Forest ensembles to effectively handle right-censored data. Therefore, this survival model allows to appropriately model data with nonlinear relationships and censoring.

\item \textbf{\acl{deepsurv}} (\acs{deepsurv})\acused{deepsurv} model~\cite{Katzman2018} (pysurvival's implementation~\cite{pysurvival}): is an improved version of the \ac{cph} model that brings in elements of deep learning to its core structure. This enhancement enables the model to better capture complex patterns while still being able to handle censored data effectively.
\end{enumerate}

    \subsection{Hyperparameters optimization}
The previous models possess distinct architectures. Consequently, the hyperparameter optimization process has varied based on the specific model.

For the purpose of identifying the optimal hyperparameters, the \ac{sa} dataset of each bike component was partitioned into three subsets: training (60\%), validation (20\%), and test (20\%). Multiple hyperparameter configurations were trained on the training set and assessed with the validation set and the RMSE metric to identify the most favorable combination. Once determined, the model was trained using both the training and validation sets, and predictions based on the test set were generated to asses the final accuracy of the model. Furthermore, although model validation was primarily assessed using the RMSE metric, all data subsets were also evaluated for RMSE, R\textsuperscript{2} and MAPE to assess accuracy comprehensively (Appendix~\ref{app_paper3:section_accuracy_measures}). RMSE and MAPE interpretation is very straightforward; the lower, the better. However R\textsuperscript{2} works in the opposite direction since higher values, correspond to better predictions.

For the \ac{cph} model, the sole hyperparameter subjected to optimization was the baseline estimation method (breslow, spline, or piecewise), while the remaining parameters were set to their default values. In the case of \ac{mtlr}, \ac{csf}, and \ac{deepsurv}, the hyperparameter search process was automated using the Optuna framework~\cite{Optuna2019}. A total of 200 trials were conducted for each model, employing the TPEsampler class for hyperparameter sampling and the MedianPruner class to halt unpromising combinations. For \ac{mtlr}, the optimized hyperparameters encompassed the learning rate (ranging between 1e-5 and 1e-3), initialization method (orthogonal or glotorot\_uniform), and optimizer (adam, adamaz, or sgd). The \ac{csf} optimization encompassed the number of trees (ranging between 10 and 100 in increments of 10), maximum tree depth (between 2 and 10), and minimum node size (ranging between 10 and 50 in steps of 5). In the case of \ac{deepsurv}, the optimization process considered initialization method (orthogonal or glotorot\_uniform), optimizer (sgd or adam), learning rate (ranging between 1e-5 and 1e-2), number of epochs (ranging between 50 and 500), L2 regularization (ranging between 0 and 1e-2), and the inclusion of batch normalization or dropout with a value of 0.5 (True or False for each one).

\section{Results}
\label{sec:failure_forecasting_results}

\subsection[MOs analysis]{\acp{mo} analysis}
\label{subsec:MOs_analysis}

    Before generating survival predictions for bike components, an analysis of the failure patterns within the maintenance dataset was conducted. Given the low frequencies of most \ac{mo} types (Appendix~\ref{app_paper3:section_data_exploration}) and the criteria outlined in Section~\ref{section_mos_processing}, only a few subcategories were deemed suitable for developing prediction models. Ultimately, three repair types were identified as the final targets for the predictions: brake pads, wheel spokes, and chains. Selected bike parts data was converted into \ac{mo} units (Table~\ref{tab:MO_units_censoring_vs_model}, Appendix~\ref{app_paper3:section_maintenance_operations_analysis}), and their corresponding covariate values were integrated as described in Section~\ref{section_mos_processing}.

    The analysis of failure dynamics for the three bike parts revealed distinct patterns between \acp{m-b} and \acp{e-b} (Figure~\ref{fig:durability_analysis}). Specifically, for brake pads and wheel spokes, \acp{e-b} generally exhibit a significantly higher number of repairs per bike compared to their mechanical counterparts. Interestingly, \acp{m-b} display very few wheel spoke repairs. Also, notable differences in the survival times of the three bike parts have been found. \Ac{e-b} brake pads have shorter survival periods compared to their mechanical counterparts, and wheel spokes display similar trends. However, the scarcity of \acp{mo} for \acp{m-b} makes direct comparisons of survival distributions impractical. Finally, chain repairs present the widest range of survival times for both \acp{m-b} and \acp{e-b}, with \ac{m-b} chains having longer survival times when compared to \ac{e-b} chains.

    \begin{figure}[ht!]
        \centering
        \includegraphics[width=1\textwidth]{0_plots/paper3/main/panel_3.pdf}
        \caption[Durability analysis of bike components]{
            \textbf{Durability analysis of bike components.} This figure presents the distribution of the survival time for the uncensored \ac{mo} units across various bike components, plotted against their cumulative traveled distance.
        }
        \label{fig:durability_analysis}
    \end{figure}


    When exploring the potential relation with bike usage, cumulative distance emerged as the primary covariate that significantly distinguishes between bike models (Figure~\ref{fig:durability_analysis}). The \ac{m-b} brake pads demonstrate considerably longer durability compared to their \ac{e-b} counterparts when covering equivalent distances, while in the case of chains, the behavior is the opposite. Furthermore, although it appears that the limited number of uncensored \ac{m-b} \ac{mo} units for wheel spokes outlast those of the \acp{e-b}, the scarcity of the former prevents a meaningful comparison of their behaviors. Also, differences have been observed in terms of the average number of reparations for other \ac{mo} subcategories during the target \ac{mo} units (Appendix~\ref{app_paper3:section_maintenance_operations_analysis}). Thus, \acp{m-b} and \acp{e-b} bike parts exhibit distinct breakdown dynamics, which can be attributed to variations in how bikes are used based on their specific model.

    Before model training, \ac{mo} units distributions according to censoring and bike model have been examined (Figure~\ref{fig:censoring}, Table~\ref{tab:MO_units_censoring_vs_model}). Uncensored data for both brake pads and wheel spokes show similar trends, representing in both cases approximately 84\% of the dataset. 90\% of the uncensored units have survival times of up to 200 days for brake pads and 150 days for wheel spokes, with some units lasting as long as approximately 800 days. However, there are two significant differences between both. First, the ratio of right-censored to uncensored units for brake pads remains fairly consistent until the 550-day mark. In contrast, for wheel spokes, the proportion of uncensored units steadily declines, virtually disappearing after 650 days. Secondly, uncensored units for \acp{m-b} become the majority for brake pads after 200 days, whereas there is an almost complete absence of such units for \ac{m-b} wheel spokes. The decreasing number of \ac{mo} units over time and the reduction in uncensored units suggest that modeling long-lasting \ac{mo} units for both bike components could pose challenges. Moreover, the specific scarcity of uncensored \ac{mo} units for \ac{m-b} wheel spokes in brake pads could introduce additional complexities in the modeling.

    \begin{figure}[ht!]
        \centering
        \includegraphics[width=1\textwidth]{0_plots/paper3/main/panel_4.png}
        \caption[Censoring distributions of bike components]{
            \textbf{Censoring distributions of bike components.} The graph features a black line representing the count of \ac{mo} units over time. Accompanying this, colored regions illustrate the relative fractions of each \ac{mo} unit type within the total count.
        }
        \label{fig:censoring}
    \end{figure}


    \begin{table}[htbp!]
        \centering
        \footnotesize
        \caption[Counts and percentages of MO units for the three chosen bike parts]{
            \textbf{Counts and percentages of \ac{mo} units for the three chosen bike parts.}
        }
        \label{tab:MO_units_censoring_vs_model}
        \begin{tabular}{lccc}
            \toprule
                            & {M-bike}        & {E-bike}         & {Total}         \\
            \midrule
            \textbf{Brake pad MO units}                                           \\
            Uncensored     & 8,777 (19.4\%)  & 29,655 (65.5\%)  & 38,432 (84.9\%) \\
            Right-censored & 3,994 (8.8\%)   & 2,841 (6.3\%)    & 6,835 (15.1\%)  \\
            Total          & 12,771 (28.2\%) & 32,496 (71.2\%)  & 45,267 (100\%)  \\ \\
        
            \textbf{Wheel spokes MO units}                                        \\
            Uncensored     & 31 (0.2\%)      & 14,177 (83.6\%)  & 14,208 (83.8\%) \\
            Right-censored & 353 (2.1\%)     & 2,386 (14.1\%)   & 2,739  (16.2\%) \\
            Total          & 384  (2.3\%)    & 16,563  (97.7\%) & 16,947 (100\%)  \\ \\
        
            \textbf{Chain MO units}                                               \\
            Uncensored     & 4,344 (36.3\%)  & 1,679 (14\%)     & 6,023 (50.4\%)  \\
            Right-censored & 3,875 (32.4\%)  & 2,056 (17.2\%)   & 5,931 (49.6\%)  \\
            Total          & 8,219 (68.8\%)  & 3,735 (31.2\%)   & 11,954 (100\%)  \\
            \bottomrule
        \end{tabular}
        \end{table}

    Chain maintenance exhibits distinct patterns. The number of censored units for chains almost exceeds that of the uncensored units, indicating that chains generally have longer lifespans compared to the other two bike components. Moreover, the failure rate for chains is notably different; only 30\% of chains fail before reaching the 200-day mark. Given the robust nature of chains, the proportion of uncensored units is low at the start of the timeline. Additionally, due to the rarity of instances surviving for extended periods, this proportion remains low past the 500-day mark as well. Considering the lower overall number of \ac{mo} units for chains, coupled with the relatively smaller number of uncensored units and their reduced proportion at both the beginning and end of the timeline, modeling the longevity of chains could be more challenging compared to other parts.

\subsection{Predictions accuracy}

    Upon confirming data distributions in the training and test sets are aligned (Appendix~\ref{app_paper3:section_model_training}), survival models were trained and their performance was evaluated predicting the uncensored \ac{mo} units from both datasets (Appendix~\ref{app_paper3:section_model_training}, Table~\ref{tab:prediction_results_test}). The decision to exclude the right-censored data was made to prevent comparisons of predictions with the time duration of \ac{mo} units that do not conclude with a repair.

    \begin{table}[htbp!]
        \centering
        \footnotesize
        \caption[Prediction accuracy on the test dataset]{
            \textbf{Predictions accuracy on the test dataset.} 
            Forecasts are computed on the test dataset using exclusively uncensored \ac{mo} units.
        }
        \label{tab:prediction_results_test}
        \begin{tabular}{llccc}
            \toprule
                                & Models   & Brake pads     & Wheel spokes   & Chains \\
            \midrule
            RMSE                & CPH      & 66.27          & 91.07          & 100.64 \\
                                & MTLR     & 38.51          & 26.74          & 64.73 \\
                                & CSF      & 37.45          & 41.28          & 108.29 \\
                                & DeepSurv & \textbf{28.75} & \textbf{18.83} & \textbf{43.62} \\
            \midrule
            R\textsuperscript{2} & CPH      & 0.59           & -0.29          & 0.60 \\
                                 & MTLR     & 0.86           & 0.89           & 0.84 \\
                                 & CSF      & 0.87           & 0.73           & 0.55 \\
                                 & DeepSurv & \textbf{0.92}  & \textbf{0.94}  & \textbf{0.93} \\
            \midrule
            MAPE                 & CPH      & 59.65          & 68.96          & 197.59 \\
                                 & MTLR     & 45.31          & 33.82          & 69.40 \\
                                 & CSF      & 89.60          & 105.90         & 227.10 \\
                                 & DeepSurv & \textbf{28.78} & \textbf{18.42} & \textbf{33.02} \\
            \bottomrule
        \end{tabular}
        \end{table}

    As expected, across nearly all models, predictions made on a combination of the training and validation sets exhibit higher accuracy than those generated on the test set. This discrepancy arises because the training and the hyper-parameter optimization process has occurred on the training and validation sets, while the test set is strictly reserved for evaluating the final performance. Moreover, the close similarity between both sets' accuracies implies that the hyper-parameter optimization has been successful enough to grant the capacity for prediction generalization.

    Classical \ac{cph} models have exhibited the poorest accuracy. This lesser performance is likely due to the their reliance on linear equations, which is inadequate for capturing the nonlinear failure dynamics of the bike components. To address this issue, \ac{mtlr} and \ac{csf} models were employed, demonstrating in both cases higher accuracies. However, an exception was noted with the \ac{csf} model's performance for chains. This outcome arises from the fact that \ac{csf} models were unable to generate predictions within the entire range of the training data, and thus, they were unable to adequately learn from the training set. In contrast, \ac{mtlr} models clearly outperformed \ac{cph} models, indicating that the adapted logistic regressions for survival data are more effective than \ac{cph} models.

    \ac{deepsurv} models were utilized to improve the modeling of nonlinear data, achieving the best results. Compared to \ac{mtlr} models, \ac{deepsurv} models demonstrated a reduction in RMSE values by 25 to 33\%. Also, they attained MAPE metrics below 33\% and R\textsuperscript{2} values exceeding 0.92. However, it is important to note that the accuracy of predictions for chains remains lower than those for brake pads and wheel spokes. This disparity might stem from two previously mentioned factors: the relatively small number of chain \ac{mo} units and the limited availability of uncensored data for \ac{mo} units with survival periods of either less than 200 days or more than 600 days.

\subsection{Predictions analysis}

    In the \ac{pm} field, predictions must be as close as possible to the failure date and, ideally, these predictions should pertain to dates preceding the failure events, as it facilitates preventive maintenance and avoids the costs associated with late component replacements. In this line, an exploration of the predictions has been performed by considering right-censored and uncensored units separately.

    Uncensored data predictions present a high level of accuracy. This assertion is further substantiated by a comparative analysis of the mean and standard deviations pertaining to the actual and predicted lifespans (Table~\ref{tab:actual_vs_pred_descriptive}). However, right-censored \ac{mo} units exhibit more substantial deviations in these statistics. Also, as previously stated, chain predictions exhibit the highest predictive errors, which can be attributed to its lower number of \ac{mo} units and the significantly higher proportion of right-censored units.

    \begin{table}[htbp!]
        \centering
        \footnotesize
        \caption[Descriptive statistics of actual and predicted survival times]{
            \textbf{Descriptive statistics of actual and predicted survival times.} 
            The means and standard deviations (in parenthesis) for the actual and predicted survival times (in days) are displayed.
        }
        \label{tab:actual_vs_pred_descriptive}
        \begin{tabular}{@{}lccc@{}}
            \toprule
                                & {All}           & {Uncensored}    & {Right-censored} \\
            \midrule
            \textbf{Brake pads MO units}                                              \\
            Actual survival    & 88.21 (96.95)   & 87.52 (97.83)   & 104.80 (103.81)  \\
            Predicted survival & 88.04 (91.19)   & 85.92 (91.86)   & 110.56 (92.42)   \\  \\
        
            \textbf{Wheel spokes MO units}                                            \\
            Actual survival    & 77.05 (88.58)   & 63.33 (71.73)   & 155.54 (126.64)  \\
            Predicted survival & 79.30 (91.62)   & 62.25 (69.58)   & 175.36 (130.84)  \\  \\
        
            \textbf{Chains MO units}                                                  \\
            Actual survival    & 258.70 (151.29) & 278.67 (147.79) & 230.81 (155.61)  \\
            Predicted survival & 283.02 (146.51) & 286.86 (145.69) & 272.73 (152.8)   \\
            \bottomrule
        \end{tabular}
        \end{table}

    Uncensored \ac{mo} units exhibit strong lineal relationships between the predicted and actual survival times, as evidenced by Pearson correlation coefficients of 0.97 for brake pads, 0.97 for wheel spokes, and 0.96 for chains (Figure~\ref{fig:pred_exploration}). Similarly, the predicted-to-actual survival time ratio for the three bike components is notably accurate (Table~\ref{tab:pred_actual_statistics}), with a mean centered around 1, and approximately 80\% of predictions exhibiting an error of less than $\pm$25\% compared to the actual values. In terms of the absolute difference, in the brake pads and wheel spokes 80\% of the \ac{mo} units have an error below 20 units, while for chains, it's 50\%. Therefore, it can be confidently concluded that the predictions are remarkably close to the actual failure dates.

        \begin{figure}[ht!]
            \centering
            \includegraphics[width=0.9\textwidth]{0_plots/paper3/main/panel_5.png}
            \caption[Comparison between actual and predicted survival times of bike components]{
                \textbf{Comparison between actual and predicted survival times of bike components.} Each plot displays instances predicted beyond the actual survival time in grey, and those predicted before in light blue. The charts also include percentages indicating the proportion of instances predicted to fail either before or after their actual failure dates.
            }
            \label{fig:pred_exploration}
        \end{figure}

        \begin{sidewaystable}[htbp!]
            \centering
            \footnotesize
            \caption[Descriptive statistics of prediction errors of uncensored MO units]{
                \textbf{Descriptive statistics of prediction errors of uncensored \ac{mo} units.}
            }
            \label{tab:pred_actual_statistics}
            \begin{tabular}{
                lcccccccccccccccc
            }
            \toprule
                            & {Mean} & {Std} & {Min} & {10\%} & {20\%} & {30\%} & {40\%} & {50\%} & {60\%} & {70\%} & {75\%} & {80\%} & {90\%} & {95\%} & {99\%} & {Max}    \\
            \midrule
            \multicolumn{17}{l}{\textbf{Predicted to actual ratio}} \\
            Brake pads   & 1.09  & 0.76  & 0.25 & 0.76 & 0.84 & 0.89  & 0.94  & 0.99  & 1.04  & 1.11  & 1.14  & 1.19  & 1.34  & 1.57  & 3.38   & 40.51  \\
            Wheel spokes & 1.03  & 0.38  & 0.00 & 0.76 & 0.84 & 0.89  & 0.93  & 0.97  & 1.02  & 1.08  & 1.12  & 1.16  & 1.28  & 1.44  & 2.20   & 13.70  \\
            Chains       & 1.20  & 2.16  & 0.00 & 0.88 & 0.92 & 0.95  & 0.98  & 1.01  & 1.04  & 1.09  & 1.11  & 1.15  & 1.30  & 1.57  & 4.57   & 112.66 \\
            \addlinespace
            \multicolumn{17}{l}{\textbf{Predicted minus actual}} \\
            Brake pads   & 12.86 & 17.84 & 0.00 & 1.07 & 2.26 & 3.61  & 5.07  & 6.85  & 9.23  & 12.81 & 15.40 & 18.84 & 31.49 & 45.59 & 86.26  & 340.79 \\
            Wheel spokes & 8.75  & 12.42 & 0.00 & 0.74 & 1.54 & 2.44  & 3.52  & 4.83  & 6.52  & 8.71  & 10.30 & 12.37 & 20.37 & 30.52 & 61.60  & 223.53 \\
            Chains       & 28.05 & 25.81 & 0.01 & 3.76 & 7.66 & 11.78 & 16.28 & 21.13 & 26.69 & 34.40 & 38.63 & 44.33 & 61.66 & 77.47 & 115.11 & 334.97 \\
            \bottomrule
            \end{tabular}  
        \end{sidewaystable}

    Unlike the uncensored units, which provide data on when bike parts fail, the right-censored units lack information regarding the specific timing of these failures. For this reason, the predicted days are expected to have higher values than the actual days, and therefore the predicted/actual ratio values should tend values higher than one. In this case, indeed, the majority of predictions generated for the right-censored \ac{mo} units were higher than the corresponding actual values (Figure~\ref{fig:pred_exploration}).

\subsection{Models interpretability}

    To assess how model inputs might affect the chances of models predicting a reparation, the game theoretic approach \ac{shap} has been applied to the \ac{deepsurv} models of the three bike parts (Figure~\ref{fig:shap}). 
    Additionally, for a comprehensive view of the average effect of each input feature on the models’ output, partial dependence plots are provided in Appendix~\ref{app_paper3:section_models_interpretability}.
    Overall, the most influential features are the ones related to the bike usage and the bike model.

    \begin{figure}[ht!]
        \centering
        \includegraphics[width=0.9\textwidth]{0_plots/paper3/main/panel_6.png}
        \caption[SHAP values for survival time predictions]{
            \textbf{\ac{shap} values for survival time predictions.} \ac{shap} values were calculated using 5,000 samples from the training set and the KernelExplainer class. This explainer was created with 100 background representative samples from the training set, which were obtained using k-means. Features are ordered according to their mean absolute \ac{shap} values in descending order.
        }
        \label{fig:shap}
    \end{figure}

    \ac{mo} units cumulative distance emerged as the most influential factor in predicting the lifespans of the three bike parts. Larger distances are strongly associated with longer predicted survival times. This connection might seem counterintuitive, as one would typically expect higher distances to lead to faster breakdowns. However, because the models' aim is to predict survival times, direct survival information is not included as an input. Consequently, due to the strong correlation between survival times and cumulative distances, cumulative distance is used as a baseline for predicting the lifespan of bike parts.

    Average mean speed has also emerged as one of the most critical features for all bike parts. Higher values are closely associated with shorter survival times, indicating that increased stress on the bike parts leads to a reduction in their lifespan.

    In Figures~\ref{fig:trips_characteristics}~and~\ref{fig:elevation_maps}, it was noted that the bike model plays an important role in \ac{bss} mobility dynamics. As a result, the bike model emerges as the second most impactful feature. Specifically, electric mobility is consistently linked with shorter survival times. However, it is worth noting that the impact of the bike model on wheel spokes is relatively less pronounced, which can be attributed to the fewer number of \ac{mo} units for \acp{m-b}.

    While weather-related variables have a lesser impact, they remain significant. Higher mean daily temperatures and mean daily atmospheric pressure values were associated with reduced survival times for brake pads and wheel spokes. In essence, hotter and drier weather conditions tended to decrease the survival time of these bike components. However, this effect is less evident in the case of chains. This discrepancy can be attributed to their extended survival times, leading to a wider range of weather conditions being encompassed within the average values. Finally, the counts of repairs for other \acp{mo} have been found to be the variables with least impact in the predictions.
        
\section{Discussion}
\label{sec:disc_maintenance}

Factors related to mobility, such as trip distance, not only differentiate electric from mechanical mobility (Chapter~\ref{ch:bicing_mobility_patterns}) but were also found to significantly influence the wear and tear on bike components. Leveraging these insights, we present a novel \ac{pm} system for a \ac{bss}, which marks a significant step towards replacing corrective maintenance strategies. The system has been designed to predict maintenance needs for three essential bike components, delivering satisfactory forecasting results through the application of \ac{dl} survival models. The design we introduce enables operation across an entire bike fleet, distinguishing it from previous research in bike \ac{pm} that primarily focused on individual bikes~\cite{MountainBike2019, Matkovic2021}. Moreover, our system enhances upon key aspects of an earlier \ac{pm} solution developed for Oslo's \ac{bss}~\cite{Predictivemaintenancebike2018}. First, our datasets are considerably more comprehensive. Second, we treat different repair typologies independently, acknowledging their distinct breakdown dynamics.  Lastly, to prevent potentially discriminatory conclusions that could impact BSS pricing strategies, we have deliberately omitted user information such as gender and age from our modeling. Additionally, through the application of a game-theoretic interpretability approach, we verify that the model's predictions are consistent with the observed failure dynamics. Notably, the most influential factors in the generation of the predictions are the cumulative distance and whether the bike is mechanical or electric.

The \ac{pm} results presented here, together with the mobility patterns documented in Chapter~\ref{ch:bicing_mobility_patterns}, hold significant promise for impacting society, especially in terms of enhancing the sustainability of urban mobility. Understanding \ac{bss} users' preferences for mechanical and electric mobility is essential for improving the decision-support systems that facilitate fleet rebalancing. Enhancing these systems not only assists \ac{bss} managers but also promotes \ac{bss} mobility by improving the user's experience, ensuring the availability of the preferred mode of transport when needed. On the other hand, the \ac{pm} system we presented, which can predict the breakdowns of three key bike components, marks an initial step towards developing a more comprehensive system that considers more bike components. Implementing such holistic systems could significantly enhance urban mobility sustainability by reducing the environmental footprint and operational costs of \ac{bss} through more efficient resource utilization. Beyond improving the scheduling of maintenance activities, these systems also facilitate the application of bike rebalancing strategies to extend the lifetime of bike parts. For instance, by strategically relocating bikes to less demanding routes or stations, we can minimize stress on critical components, thereby prolonging their usability when needed. Furthermore, the \ac{pm} systems would enhance the user's experience by ensuring a more reliable and available bike fleet. However, while this study demonstrates the feasibility of successfully deploying a \ac{pm} system for \ac{bss}, employing interpretability tools to build trust in the accuracy of these systems among \ac{bss} managers is crucial for their broader adoption. 

Although our data was of high quality, several limitations may have influenced our results. Because we lacked precise bike route information, we had to estimate trip distances, which may have impacted aspects of the \ac{pm} modeling—particularly for covariates such as cumulative distance. Similarly, we simplified elevation changes, potentially missing irregularities in the routes between stations that could affect wear-related variables. In addition, maintenance records indicated when components were replaced rather than when they actually failed. Finally, imbalances in the data due to censoring could have affected the predictive performance of our survival models.

Finally, building on the findings of this study, future research could extend the \ac{pm} framework to include more components, developing a more comprehensive system. A crucial step to validate this approach would involve the real-world deployment of these models, providing practical evidence of their impacts.
